\begin{abstract}
Although large-scale unsupervised language models (LMs) acquire extensive world knowledge and display certain reasoning capabilities, achieving precise behavioral control over these models remains challenging due to the absence of supervision during training. Approaches designed to enhance model steerability typically involve collecting human judgments comparing outputs, then adapting the unsupervised LM to reflect these preferences, most often via reinforcement learning from human feedback (RLHF). Nevertheless, RLHF entails a multi-stage and frequently unstable process: initially constructing a reward model that encodes human preferences, followed by reinforcement learning-based fine-tuning of the original LM to optimize the modeled reward while preserving the base model's features. In this work, we present a novel way to parameterize the reward model within RLHF, which makes it possible to analytically derive the corresponding optimal policy and, consequently, to address the conventional RLHF task using only a standard classification objective. Our method, termed \textit{Direct Preference Optimization} (DPO), is robust, efficient, and computationally economical; it obviates the need for sampling from the LM during adaptation or extensive hyperparameter tuning. Empirical results demonstrate that DPO matches or surpasses prior techniques in tuning LMs to match human preferences. Remarkably, DPO surpasses RLHF methods based on PPO for sentiment control, while equaling or outperforming them on tasks like summarization and single-turn conversation—all with greater simplicity and ease of training.
\end{abstract}

\section{Introduction}
Large-scale unsupervised language models (LMs) trained on massive corpora have demonstrated unexpected and impressive abilities~\citep{chowdhery2022palm, brown2020language, touvron2023llama,bubeck2023sparks}. Yet, these models learn from human-generated content characterized by a diverse array of objectives, priorities, and competencies. Some of these objectives and competencies are not always appropriate for imitation; for instance, an AI coding assistant ought to \textit{recognize} common programming errors to rectify them, but in the context of code generation, we want to promote the (sometimes scarce) examples of expert-level coding present in its dataset. Analogously, we may require the model to be \textit{informed} about prevalent misconceptions—such as those held by 50\% of individuals—but it is undesirable for the LM to affirm the misconception in half of its relevant outputs. Therefore, discerning and privileging the model's \emph{target behaviors and outputs} from the broader spectrum of its \textit{acquired knowledge and competencies} is vital for achieving safe, effective, and manageable AI systems~\citep{ouyang2022training}. While most existing approaches steer LMs towards human-aligned behaviors using reinforcement learning (RL), we will demonstrate that the RL objectives implemented by standard methodologies can be precisely addressed by a straightforward binary cross-entropy loss, streamlining the entire preference optimization process.

At a conceptual level, mainstream techniques impart desirable behavioral traits to a language model by leveraging carefully assembled datasets of human preference judgements that delineate safe and useful actions. This phase of preference alignment follows an initial period of extensive, unsupervised pre-training on text corpora. The simplest method for aligning models with preferences is through supervised fine-tuning on exemplary human responses, but the most effective contemporary approaches are based on reinforcement learning from human (or AI) feedback (RLHF/RLAIF; \citep{christiano2017deep,bai2022constitutional}). RLHF approaches involve first constructing a reward model via human-labeled preference data, then optimizing the LM policy with RL to maximize predicted rewards while constraining deviation from the base model. Although RLHF results in models with outstanding dialog and programming performance, the procedure introduces considerable complexity beyond standard supervised learning; it requires training several language models and repeatedly sampling outputs during RL training loops, thereby entailing substantial computational expenditure.

In this work, we introduce a method for optimizing language models to align with human preferences directly, foregoing the need for explicit reward modeling or reinforcement learning. Our method, termed \textit{{Direct Preference Optimization} (DPO)}, provides a mechanism for implicitly maximizing the same reward-plus-KL objective common to RLHF approaches, yet remains more straightforward to implement and train. At its core, DPO modifies the ratio of log probabilities between preferred and non-preferred outputs, but crucially, it integrates a dynamic, sample-specific importance weighting. This weighting acts to avert the model degradation that tends to arise when using a naive probability ratio-based approach. Like prior works, DPO is grounded in a formal preference framework—such as the Bradley-Terry model (\cite{bradley1952rankanalysis})—that quantifies the agreement between a reward function and observed human preference data. Distinct from existing approaches that employ this preference structure to define a loss for learning a reward model, and then separately optimize a policy against this learned reward, DPO reformulates the loss to depend directly on the policy through a variable transformation. Provided with a dataset of human preference judgments over model outputs, DPO is able to train the model with a simple binary cross entropy loss, effectively yielding the policy optimal with respect to an implicit reward tailored to the observed preferences.

The principal innovation of this paper is the Direct Preference Optimization (DPO) framework, which enables preference-based language model training without relying on reinforcement learning. Through empirical studies, we demonstrate that DPO matches or outperforms state-of-the-art methods, such as RLHF with PPO, across various tasks—including sentiment control, summarization, and conversational modeling—on language models with sizes up to 6B parameters.

\section{Related Work}

Self-supervised language models, as their scale grows, acquire the capacity to handle certain tasks without explicit training examples (zero-shot) \citep{radford2019language}, or when provided only with a handful of demonstrations (few-shot) \citep{gpt3,megatron,chowdhery2022palm}. Nevertheless, the effectiveness of these models on subsequent tasks and their alignment with user expectations is greatly enhanced through additional training on corpora composed of instructional prompts and human-generated completions \citep{mishra-etal-2022-cross,sanh2022multitask,chung2022scaling,thoppilan2022lamda}. This method, termed `instruction-tuning,’ facilitates broader generalization by LLMs to tasks not explicitly seen during fine-tuning and consistently improves overall user interaction quality \citep{chung2022scaling}. In spite of its advantages, gathering \textit{relative} judgments from humans about which responses are better is generally more scalable than acquiring full expert-labeled demonstrations. Consequently, recent research has pursued further fine-tuning of LLMs using collections of human preference data, leading to advances in translation \citep{kreutzer-etal-2018-reliability}, summarization \citep{stiennon2022learning,ziegler2020finetuning}, narrative generation \citep{ziegler2020finetuning}, and obedience to instructions \citep{ouyang2022training,ramamurthy2023is}. These approaches initially involve training a neural reward model to reflect preferences using models such as the Bradley-Terry preference framework \citep{bradley1952rankanalysis}, followed by optimizing the language model itself to maximize this reward function, typically employing reinforcement learning techniques such as REINFORCE \citep{williams1992reinforce}, proximal policy optimization (PPO; \cite{schulman2017proximal}), or related strategies \citep{ramamurthy2023is}. An adjacent research area makes use of instruction-tuned LLMs further enhanced with human feedback, in order to synthesize additional preference data focused on qualities like safety or harmlessness \citep{bai2022constitutional}. Here, human supervision is often indirect, supplied as rubric-based textual guidance for LLM-generated annotations. Collectively, these developments merge two research streams: the use of reinforcement learning to train language models for various objectives~\citep{Ranzato2015SequenceLT,paulus2018a,wu2018learning}, and generic methodologies for learning from human preferences \citep{christiano2017deep,kupcsik2018learning}. Despite the promise of leveraging comparative human preferences, reinforcement learning-based fine-tuning of large language models still presents significant practical barriers; the present work introduces a theoretically grounded alternative for optimizing with relative preferences that avoids reliance on RL.

Research on learning policies from preferences, outside of language domains, has taken place within both bandit and reinforcement learning frameworks, leading to the development of several methodologies. In contextual bandit settings, learning with preferences or rankings over actions—as opposed to explicit rewards—is referred to as the contextual dueling bandit (CDB) problem \citep{yue2012karmed,dudik2015contextual}. Lacking access to absolute rewards, theoretical treatments of CDBs rely on the concept of a \textit{von Neumann winner}, that is, a policy guaranteed to achieve at least a 50\% expected win rate against \textit{every} competing policy \citep{dudik2015contextual}. Nevertheless, the CDB framework assumes that preference information is received in an online manner, whereas in the context of learning from human preferences, practitioners usually depend on a fixed, offline dataset comprised of pairs of actions annotated with preferences \citep{yan2022human}. Relatedly, \textit{preference-based RL} (PbRL) learns from binary feedback derived from an \textit{unknown} scoring function, rather than relying on traditional reward signals \citep{BusaFekete2014,ruiz2023dueling}. A range of PbRL algorithms has been proposed, some capable of utilizing off-policy preference data; these approaches commonly require first modeling the underlying scoring function (that is, constructing the reward model) and then optimizing policies with respect to this estimate \citep{jain2013learning,BusaFekete2014,christiano2017deep,sadigh2017active,kupcsik2018learning}. By contrast, we introduce a one-step policy learning method that directly maximizes compliance with preferences, foregoing the intermediate reward modeling stage.

\section{Preliminaries}\label{section:prelims}

We provide an overview of the RLHF procedure as detailed in \citeauthor{ziegler2020finetuning} and further employed in subsequent works \citep{stiennon2022learning, bai2022training, ouyang2022training}. The process generally consists of three main stages: (1) supervised fine-tuning (SFT), (2) collecting preference data and fitting a reward function, and (3) optimizing the model through reinforcement learning.

\textbf{SFT}: The initial step in RLHF involves adapting a large, pre-trained language model to the intended application (such as dialogue or summarization) via supervised learning on curated, high-quality examples, producing the SFT model $\pi^\text{SFT}$.

\textbf{Reward Modelling Phase}: Next, the SFT model generates answer pairs $(y_1, y_2)\sim \pi^\text{SFT}(y \mid x)$ in response to given prompts $x$. Human annotators then compare these outputs and indicate their preferred response; this is represented as $y_w\succ y_l \mid x$, where $y_w$ is the chosen completion and $y_l$ is the less favored alternative from the set $(y_1, y_2)$. The actual preferences are presumed to stem from an unobserved reward function $r^*(y, x)$. Various methods have been proposed for modeling these preferences, with the Bradley-Terry (BT) model \cite{bradley1952rankanalysis} being particularly prevalent (though the approach readily accommodates more sophisticated models, such as Plackett-Luce \citep{plackett1975analysis, luce2012individual}, when ranked lists of completions are available). In the BT framework, the distribution over human preferences $p^*$ can be expressed as follows:

Given a fixed dataset of pairwise preferences, $\mathcal{D}=\bigl\{x^{(i)}, y_w^{(i)}, y_l^{(i)}\bigr\}_{i=1}^N$, sampled according to $p^*$, one can represent the reward function as $r_{\phi}(x, y)$ and infer its parameters via maximum likelihood estimation. By reframing this as a binary classification challenge, the model is trained to minimize the negative log-likelihood objective:

\begin{equation}\label{eq:reward_model}
    \mathcal{L}_R(r_{\phi}, \mathcal{D}) = -\mathbb{E}_{(x, y_w, y_l)\sim \mathcal{D}}\bigl[\log \sigma(r_{\phi}(x, y_w)- r_{\phi}(x, y_l))\bigr]
\end{equation}

with $\sigma$ denoting the logistic sigmoid function. When applying this to language models (LMs), it is standard to initialize $r_{\phi}(x, y)$ from the SFT model $\pi^\text{SFT}(y \mid x)$, appending a linear output head to the final transformer representation, producing a scalar reward estimate as suggested by \cite{ziegler2020finetuning}. To promote stability, the reward outputs are typically normalized such that $\mathbb{E}_{x,y\sim \mathcal{D}}\left[r_\phi(x, y)\right] = 0$ for each input $x$.

\textbf{RL Fine-Tuning Phase}: In the reinforcement learning step, the previously learned reward model furnishes feedback to guide the language model’s updates. As in prior research~\citep{jaques2017sequence, jaques2020human}, the optimization task is formulated as

\begin{equation}\label{eq:RL}
\max_{\pi_{\theta}}  \mathbb{E}_{x\sim \mathcal{D}, y\sim \pi_{\theta}(y \mid x)}\bigl[r_{\phi}(x, y)\bigr] - \beta\mathbb{D}_{\textrm{KL}}\bigl[\pi_{\theta}(y\mid x)\mid \mid \pi_\text{ref}(y\mid x)\bigr],
\end{equation}

where the hyperparameter $\beta$ governs how far the learned policy $\pi_\theta$ is permitted to diverge from a reference distribution, namely $\pi_\text{ref}$—usually the SFT policy $\pi^\text{SFT}$. In practice, $\pi_\theta$ is initialized with $\pi^\text{SFT}$. Imposing this KL penalty is vital: it restricts excessive divergence from regions of the distribution where the reward model remains valid, preserves generative variety, and counters tendencies for the model to converge on narrow sets of high-reward sequences (“mode collapse”). Since language generation involves discrete choices, gradients cannot be propagated directly; reinforcement learning methods are therefore used to optimize this objective. A standard approach \citep{ziegler2020finetuning, stiennon2022learning, bai2022training, ouyang2022training} is to define the scalar reward as ${r(x, y) = r_{\phi}(x, y) -\beta (\log \pi_{\theta}(y\mid x) - \log \pi_\text{ref}(y\mid x))}$, and maximize it via the PPO algorithm \cite{schulman2017proximal}.

\section{Direct Preference Optimization}\label{sec:DPO}

Recognizing the complexity that arises when deploying reinforcement learning algorithms at scale—such as when fine-tuning LMs—our objective is to design a straightforward policy optimization procedure that can make direct use of preference data. Contrary to conventional RLHF pipelines, which train an explicit reward model and then optimize the model using RL, our approach utilizes a parameterization of the reward such that the optimal policy can be characterized analytically, making the RL step unnecessary. We show that an analytical mapping exists between reward models and their corresponding optimal policies; by re-expressing the learning objective from reward-space to policy-space, we eliminate the need for separate, explicit reward model fitting. Instead, our approach directly optimizes the policy under preference data assumptions, like those used in the Bradley-Terry framework. Thus, a single policy network fulfills both the generation and implicit reward roles.

\textbf{Derivation of the DPO objective.} We begin by considering the reinforcement learning objective given in Eq.~\ref{eq:RL}, as used in previous work, under a generic reward function $r$. Prior literature~\citep{peters2007reinforcement, peng2019advantage, korbak2022reinforcement, go2023aligning} establishes that the optimal solution to the reward maximization problem subject to a KL constraint, as expressed in Eq.~\ref{eq:RL}, adopts the following form:

\begin{equation}\label{eq:op_policy}
    \pi_r(y\mid x) = \frac{1}{Z(x)}\pi_\text{ref}(y\mid x)\exp\left(\frac{1}{\beta}r(x, y)\right),
\end{equation}

with $Z(x) =\sum_{y}\pi_\text{ref}(y\mid x)\exp\left(\frac{1}{\beta}r(x, y)\right)$ serving as the normalization constant or partition function. Further details and a full derivation are available in Appendix \ref{app:derivation1}. Even when substituting an MLE-based estimate $r_{\phi}$ for the true reward function $r^*$, evaluating the partition function $Z(x)$ remains computationally demanding \citep{korbak2022reinforcement, go2023aligning}, complicating the direct application of this solution. Nevertheless, by manipulating Eq.~\ref{eq:op_policy}, we can re-express the reward in terms of its induced optimal policy $\pi_r$, the reference policy $\pi_\text{ref}$, and the generally intractable partition function $Z(\cdot)$. To do this, we apply the logarithm to both sides of Eq.~\ref{eq:op_policy} and rearrange terms to find:

\begin{equation}\label{eq:main_eq}
    r(x,y) =\beta \log \frac{\pi_r(y\mid x)}{\pi_\text{ref}(y\mid x)} + \beta \log Z(x).
\end{equation}

This reparameterization can be similarly applied to the true reward $r^*$ and its associated optimal policy $\pi^*$. Importantly, the Bradley-Terry preference model utilizes only the reward difference between two candidates, i.e., ${p^*(y_1 \succ y_2 \mid x) = \sigma(r^*(x, y_1) - r^*(x, y_2))}$. Substituting the relation from Eq.~\ref{eq:main_eq} for $r^*(x,y)$ into Eq.~\ref{eq:bradley-terry}, the partition function terms eliminate, leaving us with a preference probability that depends exclusively on $\pi^*$ and $\pi_\text{ref}$. Therefore, under the Bradley-Terry model, the optimal RLHF policy $\pi^*$ satisfies:

\begin{equation}\label{eq:objective}
    p^*(y_1\succ y_2 \mid x)=\frac{1}{1 + \exp\left(\beta \log \frac{\pi^*(y_2\mid x)}{\pi_\text{ref}(y_2\mid x)} - \beta \log \frac{\pi^*(y_1\mid x)}{\pi_\text{ref}(y_1\mid x)}\right)}
\end{equation}

For a full exposition of this result, see Appendix~\ref{app:derivation2}. Although Eq.~\ref{eq:objective} is grounded in the Bradley-Terry model, an analogous treatment yields similar formulations for broader preference models, such as the Plackett-Luce family~\citep{plackett1975analysis, luce2012individual}, as elaborated in Appendix~\ref{app:plackett_luce_models}.

Having now characterized the human preference probability in terms of the optimal policy rather than the reward, we are able to specify a maximum likelihood learning objective for a parameterized policy $\pi_\theta$. Mirroring the strategy of reward modeling (as in Eq.~\ref{eq:reward_model}), our policy training objective is formulated as:

\begin{equation}\label{eq:optimum_model}
    \mathcal{L}_\text{DPO}(\pi_{\theta}; \pi_\text{ref}) = -\mathbb{E}_{(x, y_w, y_l)\sim \mathcal{D}}\left[\log \sigma \left(\beta \log \frac{\pi_{\theta}(y_w\mid x)}{\pi_\text{ref}(y_w

In this formulation, we learn an implicit reward by adopting a different parameterization, wherein the resulting optimal policy coincides with $\pi_\theta$. Furthermore, because our approach mirrors the structure of a reparametrized Bradley-Terry model, it inherits advantageous theoretical guarantees, including consistency provided that the preference data distribution satisfies appropriate conditions \cite{bong2022generalized}. Section~\ref{sec:theory} provides a more comprehensive analysis of the theoretical properties of DPO and how it relates to existing literature.

\textbf{How does the DPO update operate?} To gain mechanistic insights into DPO, it is instructive to investigate the gradient of its loss $\mathcal{L}_\text{DPO}$. The derivative of this objective with respect to the parameter vector $\theta$ is expressed as:

\begin{multline*}\label{eq:gradient}
    \nabla_\theta \mathcal{L}_\text{DPO}(\pi_\theta;\pi_\text{ref}) = \\ -\beta\mathbb{E}_{(x, y_w, y_l) \sim \mathcal{D}} \bigg[\underbrace{\sigma(\hat{r}_\theta(x, y_l) - \hat{r}_\theta (x, y_w))}_\text{assigns greater emphasis when reward prediction errs}\bigg[\underbrace{\nabla_\theta\log \pi(y_w \mid x)}_\text{promote $y_w$ probability} - \underbrace{\nabla_\theta\log\pi(y_l \mid x)}_\text{reduce $y_l$ probability}\bigg]\bigg],
\end{multline*}

where the reward is implicitly defined as $\hat{r}_\theta(x, y) = \beta \log \frac{\pi_\theta(y \mid x)}{\pi_\text{ref}(y \mid x)}$, as determined by both the current model $\pi_\theta$ and the reference model $\pi_\text{ref}$ (see Section~\ref{sec:theory} for further discussion). Conceptually, this gradient updates the model to assign higher likelihood to favored completions $y_w$ and suppress the likelihood of less favored completions $y_l$. Crucially, each example is weighted based on the extent to which the implicit reward model $\hat{r}_\theta$ ranks the disfavored completion too high, modulated by the parameter $\beta$; in essence, the weighting reflects the degree of error in the completion ordering imposed by the reward model, adjusted by the strength of the KL regularization. Our empirical findings highlight the criticality of this weighting: omitting the weighting factor—effectively a simplified variant—leads to pathological behavior in the language model, as illustrated in Appendix Table~\ref{tab:unlikelihood_generations}.

\textbf{DPO overview.} The standard DPO framework consists of the following steps: 1) For each prompt $x$, generate completions $y_1, y_2$ from $\pi_\text{ref}(\cdot \mid x)$, and annotate them using human preference judgments to assemble an offline preference dataset $\mathcal{D} = \{x^{(i)}, y_w^{(i)}, y_l^{(i)}\}_{i=1}^N$; 2) Train the language model $\pi_\theta$ by minimizing $\mathcal{L}_\text{DPO}$ given $\pi_\text{ref}$, $\mathcal{D}$, and a chosen $\beta$. In realistic scenarios, practitioners often seek to leverage existing public preference datasets rather than generate new samples and solicit fresh human feedback. Because these datasets are generally created using $\pi^\text{SFT}$ as the sampling distribution, we set $\pi_\text{ref} = \pi^\text{SFT}$ when it is accessible. If $\pi^\text{SFT}$ is absent, we instead fit $\pi_\text{ref}$ by maximizing the likelihood over preferred completions $(x, y_w)$ via ${\pi_\text{ref} = \argmax_{\pi}\mathbb{E}_{x, y_w \sim \mathcal{D}}\left[\log \pi(y_w \mid x)\right]}$. This initialization strategy reduces the mismatch between the ideal but unavailable reference distribution and the $\pi_\text{ref}$ employed by DPO. Additional details concerning implementation choices and hyperparameter selection are provided in Appendix~\ref{app:implementation}.

\section{Theoretical Analysis of DPO} This section develops the theoretical foundations for DPO, offering further insight into the approach, presenting justification, and comparing DPO’s benefits with challenges faced by actor-critic frameworks in RLHF, such as PPO~\cite{schulman2017proximal}.

\label{sec:theory}

\subsection{Your Language Model Is Secretly a Reward Model} DPO eliminates the need for explicitly modeling rewards or running separate RL optimization for the policy, by utilizing a unified maximum likelihood framework. Specifically, the objective in Eq.~\ref{eq:main_eq} corresponds to a Bradley-Terry model where the reward function is given by $r^*(x, y) = \beta \log\frac{\pi^*_\theta(y \mid x)}{\pi_\text{ref}(y \mid x)}$, and training the parametric model $\pi_{\theta}$ is mathematically equivalent (under a suitable change of variables) to reward model optimization as shown in Eq.~\ref{eq:reward_model}. In this section, we formalize the theoretical rationale underlying this reparameterization, demonstrate that it does not reduce the generality of the reward models that can be learned, and prove that the procedure can exactly recover the optimal policy. We commence by introducing an equivalence notion for reward functions.

\begin{definition}
We define two reward functions $r(x, y)$ and $r'(x, y)$ to be equivalent if and only if ${r(x, y) - r'(x, y) = f(x)}$ for some function $f$.
\end{definition}

It is straightforward to verify that this constitutes an equivalence relation, which divides the space of reward functions into distinct equivalence classes. This allows us to formalize the following two results:

\begin{lemma}\label{lemma:same_prefrence}
Within the framework of the Plackett-Luce, and more specifically the Bradley-Terry, models for preferences, any pair of reward functions belonging to the same equivalence class generates an identical preference distribution.
\end{lemma}

\begin{lemma}\label{lemma:same_policy}
    For the constrained RL setting, any two reward functions drawn from the same equivalence class give rise to the same optimal policy.
\end{lemma}

Since the proofs are elementary, we omit them here and provide details in Appendix \ref{app:lemma1}. The first lemma highlights a commonly recognized under-identification property of the Plackett-Luce family \cite{plackett1975analysis}. Consequently, additional constraints ensuring identifiability are typically introduced in order to obtain reliable MLE solutions for Eq. \ref{eq:reward_model} \cite{bong2022generalized}. The second lemma demonstrates that every member of a given equivalence class yields the same optimal policy, which means our main objective is to recover any representative reward function from the optimal equivalence class. In Appendix~\ref{app:thm1}, we establish the following theorem:

\begin{theorem}\label{thm:main}
    Assuming mild conditions, every equivalence class of reward functions consistent with the Plackett-Luce (and specifically the Bradley-Terry) model admits a reparameterization of the form ${r(x, y) = \beta \log \frac{\pi(y\mid x)}{\pi_\text{ref}(y\mid x)}}$ for some policy model $\pi(y\mid x)$ and a given reference model $\pi_\text{ref}(y \mid x)$.
\end{theorem}

\begin{sproof}
    Let $r(x, y)$ denote a reward function associated with an optimal policy $\pi_r(y \mid x)$ as described by Eq. \ref{eq:op_policy}. We will demonstrate that there exists a reward function in the equivalence class of $r$ that can be expressed using the proposed reparameterization. Define the projection operator $f$ by  
\begin{equation}
    f(r; \pi_\text{ref}, \beta)(x, y) = r(x, y) - \beta\log\sum_{y}\pi_\text{ref}(y\mid x)\exp\left(\frac{1}{\beta}r(x, y)\right)
\end{equation}
This operator $f$ effectively normalizes $r(x, y)$ by subtracting the logarithm of the policy's partition function. As this normalization term depends solely on $x$, $f(r; \pi_\text{ref}, \beta)(x, y)$ remains within the equivalence class of $r(x, y)$. Substituting the right side of Eq.~\ref{eq:main_eq} (valid for all reward functions) in place of $r$ yields $f(r; \pi_\text{ref}, \beta)(x, y) = \beta \log \frac{\pi_r(y\mid x)}{\pi_\text{ref}(y\mid x)}$. Therefore, $f$ maps any reward function in the class to an element with the reparameterized form, ensuring no representational loss with the chosen parametrization of the reward.
\end{sproof}

An alternative interpretation of Theorem~\ref{thm:main} is that it identifies precisely which reward function, among those within an equivalence class, is chosen by the DPO reparameterization; specifically, it is the reward function that satisfies:

\begin{equation}\label{eq:lag_p}
     \sum_{y}\underbrace{\pi_\text{ref}(y\mid x)\exp\left(\frac{1}{\beta}r(x, y)\right)}_{=\pi(y\mid x)\text{, using Thm.~\ref{thm:main} reparam.}} = 1,
\end{equation}

meaning that $\pi(y\mid x)$ forms a proper probability distribution: all probabilities are non-negative and collectively sum to one. Moreover, by considering Eq.~\ref{eq:op_policy}, it becomes clear that Eq.~\ref{eq:lag_p} represents the partition function for the optimal policy that the reward function $r(x, y)$ induces. The crucial insight underlying the DPO algorithm is that, by enforcing appropriate constraints on the inherently under-constrained Plackett-Luce (and more specifically, Bradley-Terry) class of preference models, we retain the set of reward models that can be represented, while at the same time rendering the optimal policy in Eq.~\ref{eq:op_policy} explicitly computable for any prompt $x$.

\subsection{Instability of Actor-Critic Algorithms} Our framework is also applicable to the analysis of instability issues in conventional actor-critic algorithms employed for RLHF, such as PPO. By adopting the RLHF process and concentrating on the RL fine-tuning stage described in Section \ref{section:prelims}, we can relate our analysis to the control as inference framework \cite{levine2018reinforcement} for the constrained RL formulation in \ref{eq:RL}. Suppose the policy is parameterized as $\pi_{\theta}(y\mid x)$. The objective is to minimize $\mathbb{D}_{\text{KL}}[\pi_{\theta}(y|x) \mid \mid \pi^*(y\mid x)]$, where $\pi^*$, as defined in Eq.~\ref{eq:optimum_model}, denotes the optimal policy corresponding to the reward $r_{\phi}(y, x)$. Through algebraic manipulation, the resulting optimization problem becomes:

\begin{equation}\label{eq:AC}
    \max_{\pi_{\theta}}\mathbb{E}_{\pi_{\theta}(y\mid x)}\left[\underbrace{r_{\phi}(x, y) -\beta\log\sum_{y}\pi_\text{ref}(y\mid x)\exp\left(\frac{1}{\beta}r_{\phi}(x, y)\right)}_{f(r_{\phi}, \pi_\text{ref}, \beta)} - \underbrace{\beta\log\frac{\pi_{\theta}(y\mid x)}{\pi_\text{ref}(y\mid x)}}_{\text{KL}}\right]
\end{equation}

This objective aligns with that optimized in earlier studies 
\citep{ziegler2020finetuning, stiennon2022learning, bai2022training, ouyang2022training} by employing the DPO-equivalent reward formulated for the reward class $r_{\phi}$. Under this framework, the normalization component in $f(r_{\phi}, \pi_\text{ref}, \beta)$ can be viewed as the soft value function corresponding to the reference policy $\pi_\text{ref}$. Although this term does not influence the optimal solution directly, omitting it may result in a high-variance policy gradient, thus causing instability during training. This normalization can, in principle, be addressed by employing a learned value function, though this approach introduces its own optimization challenges. An alternative strategy used in prior research involves normalizing the rewards relative to a human completion baseline, which essentially serves as a one-sample Monte Carlo approximation of the normalization factor. Distinctly, the DPO reparameterization gives rise to a reward function that is free from any reliance on such baselines.

\section{Experiments}
This section provides an empirical assessment of the effectiveness of DPO for policy training directly from preference data. We begin by conducting experiments in a controlled text-generation environment to investigate: how effectively does DPO balance reward maximization with maintaining closeness to the reference policy (in terms of KL-divergence) relative to standard preference-based algorithms such as PPO? Subsequently, we extend our analysis to encompass larger-scale models and more complex RLHF tasks, including dialogue and summarization. Across these experiments, we observe that DPO, with minimal hyperparameter adjustment, often matches or exceeds the performance of robust baselines like RLHF with PPO and even outperforms selecting the best out of $N$ sampled trajectories as determined by a learned reward function. Prior to discussing these empirical findings, we outline our experimental protocol; comprehensive experimental details are available in Appendix~\ref{app:exp_details}.

\textbf{Tasks.} Our study investigates three distinct open-ended text generation tasks. Across all tasks, the learning objective is to develop a policy from a collection of preference data $\mathcal{D}=\bigl\{x^{(i)}, y_w^{(i)}, y_l^{(i)}\bigr\}_{i=1}^N$. In the case of \textbf{controlled sentiment generation}, $x$ denotes the beginning segment of a movie review sourced from the IMDb dataset \cite{maas-EtAl:2011:ACL-HLT2011}, with the requirement that the policy outputs $y$ exhibiting positive sentiment. To enable a systematic evaluation, for this scenario, preference pairs over generated texts are constructed using a pre-trained sentiment classifier, ensuring $p(\text{positive}\mid x,y_w)>p(\text{positive}\mid x,y_l)$. For SFT, the GPT-2-large model undergoes further training on the training portion of the IMDb reviews dataset until convergence (see additional details in App~\ref{app:sentiment_details}). In the \textbf{summarization} task, $x$ is a Reddit forum post, and the goal is for the policy to generate a concise summary $y$ that covers the main content of the post. Consistent with earlier research, we utilize the Reddit TL;DR summarization dataset \citep{volske-etal-2017-tl} in conjunction with human preference data provided by \citeauthor{stiennon2022learning}. Here, the SFT model is adapted through fine-tuning on human-written summaries of forum posts\footnote{\url{https://huggingface.co/CarperAI/openai_summarize_tldr_sft}} and uses the TRLX \citep{leandro_von_werra_2023_7790115} toolkit for RLHF. The corresponding human preference data was collected by \citeauthor{stiennon2022learning} from outputs of a related but not identical SFT model. Lastly, for the \textbf{single-turn dialogue} task, $x$ serves as a user prompt, which could range from scientific inquiries to requests for advice. Here, the policy is expected to generate an informative and engaging reply $y$; the dataset used is the Anthropic Helpful and Harmless dialogue corpus \citep{bai2022training}, comprising 170,000 dialogues between humans and an automated assistant. Each dialogue concludes with two responses—both generated by a large, unspecified language model—accompanied by an indicator of the human-preferred reply. In the absence of a suitable pre-trained SFT model for this domain, we create an SFT model by fine-tuning an existing language model exclusively on the responses designated as preferred.

\textbf{Evaluation.} Our evaluation framework employs two distinct methodologies. To rigorously assess each algorithm’s ability to optimize the constrained reward maximization objective, we conduct controlled sentiment generation experiments, plotting each algorithm's achieved reward against its KL-divergence relative to the reference policy; this evaluation is feasible given our access to the true reward function via a sentiment classifier. In contrast, when the actual reward function is unavailable—as is the case in practical deployment—we measure each algorithm's \textit{win rate} against a baseline policy. For this, GPT-4 serves as a stand-in for human assessment, scoring both summary quality (in the summarization setting) and response helpfulness (in single-turn dialogue tasks). Baselines consist of test set reference summaries for summarization and the dataset's preferred responses for dialogue. Although recent work suggests language models may outperform traditional metrics as evaluators \citep{Chen2023ExploringTU}, we include a human study (see Sec.~\ref{sec:human-judgments}) to support the reliability of GPT-4 for evaluation. Our findings indicate a strong alignment between GPT-4’s and human judgments, with the degree of human agreement with GPT-4 equaling or surpassing that found between human annotators.

\textbf{Methods.} Besides DPO, we investigate multiple established methods for steering language models toward human-preferred behaviors. We apply zero-shot prompting using \textbf{GPT-J} \citep{gpt-j} in the summarization context and two-shot prompting with \textbf{Pythia-2.8B} \citep{biderman2023pythia} for dialogue. Our assessment also includes the \textbf{SFT} model and a variant named \textbf{Preferred-FT}, which is trained through supervised fine-tuning on the selected completions $y_w$—originating from either the SFT model (for sentiment and summarization) or a general language model (in dialogue). An alternative pseudo-supervised strategy is \textbf{Unlikelihood}~\citep{welleck2019neural}, where the objective is to enhance the probability of $y_w$ while explicitly penalizing the likelihood of $y_l$, incorporating a tunable weight $\alpha\in[0,1]$ for the `unlikelihood' penalty. We further analyze \textbf{PPO} \citep{schulman2017proximal} with a reward signal trained on preference data, as well as \textbf{PPO-GT}, which represents an oracle scenario with access to the genuine reward in controlled sentiment generation. For the sentiment tasks, two PPO-GT implementations are used: a standard off-the-shelf version \cite{leandro_von_werra_2023_7790115} and a customized version that includes reward normalization and refined hyperparameter settings for optimal outcomes (these modifications are also applied to PPO runs using learned rewards). Lastly, we benchmark against the \textbf{Best of $N$} strategy, where $N$ outputs are sampled from the SFT model (or Preferred-FT in the dialogue context), and the top response according to a reward model trained on preference data is selected. While this approach often yields high-quality outputs and disentangles reward model quality from PPO’s optimization process, it is computationally demanding—even for modest values of $N$—as it necessitates generating $N$ candidate responses for each test input.

\subsection{How well can DPO optimize the RLHF objective?}

The standard KL-constrained reward maximization framework in most RLHF methodologies seeks a compromise between maximizing reward and constraining the policy's divergence from a reference policy. Thus, when evaluating different algorithms, it is important to consider both the level of reward obtained and the magnitude of the KL divergence; it is not always advantageous to gain higher reward at the expense of a much larger KL divergence. Figure~\ref{fig:frontier-tldr-main} illustrates the tradeoff curve between reward and KL divergence for several algorithms applied in the sentiment task. For each algorithm, we conduct multiple training sessions, varying the policy conservativeness hyperparameters across runs (using target KL values $\in\{3,6,9,12\}$ for PPO, $\beta \in \{0.05,0.1,1,5\}$, and $\alpha\in\{0.05,0.1,0.5,1\}$ for unlikelihood, as well as different random seeds for preferred-FT), yielding a total of 22 experiments. At every 100 training iterations until training completion, each resulting policy is evaluated on a set of held-out prompts, with the mean reward under the ground truth reward function and the average sequence-level KL\footnote{Specifically, this is computed as the aggregate of the per-timestep KL-divergences.} with respect to the reference policy, $\text{KL}\left(\pi\mid \mid \pi_\text{ref}\right)$, being reported. Our results indicate that DPO generates the most favorable frontier, securing superior reward outcomes while maintaining low KL divergence. This finding is significant for several reasons. First, while both DPO and PPO optimize the identical objective, DPO is markedly more efficient, with its reward/KL tradeoff surpassing that of PPO in all cases. Furthermore, DPO attains a superior tradeoff compared to PPO, \emph{even in settings where PPO is provided with access to the true rewards} (PPO-GT).

\subsection{Can DPO scale to real preference datasets?}
\label{sec:dpo-real-datasets}
We proceed by assessing the fine-tuning capabilities of DPO on tasks such as summarization and single-turn dialogue. For summarization, traditional metrics like ROUGE often exhibit weak alignment with human judgments~\citep{stiennon2022learning}, and previous studies have indicated that optimizing language models via PPO on human-labeled preferences leads to higher-quality summaries. In our experiments, we compare several approaches by generating completions using the test portion of the TL;DR summarization dataset and calculating the average frequency with which each method outperforms reference completions. To evaluate performance across different settings, we sample outputs from each model using temperatures in the range 0.0 to 1.0, presenting win rates in Figure~\ref{fig:frontier-tldr-main} (right). DPO, PPO, and Preferred-FT all use the same GPT-J SFT checkpoint\footnote{\url{https://huggingface.co/CarperAI/openai_summarize_tldr_sft}} for further training. Our analysis indicates that DPO attains a win rate of about 61\% at a sampling temperature of 0.0, surpassing PPO, which achieves roughly 57\% under its optimal temperature condition (also 0.0). Furthermore, DPO secures a higher peak win rate than the best-of-$N$ baseline. It is worth noting that we did not undertake extensive hyperparameter optimization for DPO, particularly for $\beta$, which suggests these results may not reflect the full potential of DPO. Notably, DPO demonstrates greater stability with respect to sampling temperature; in contrast, PPO's efficacy can decline to the level of the original GPT-J model at higher temperatures. Improvements from Preferred-FT over the SFT baseline were marginal. Finally, a direct human evaluation between DPO and PPO (see Section~\ref{sec:human-judgments}) shows that completions from DPO at temperature 0.25 were favored 58\% of the time over those from PPO at the same temperature.

For single-turn dialogue evaluation, we assess each approach on a subset of the Anthropic HH dataset's test split \citep{bai2022training}, restricting the interaction to a single exchange between human and assistant. In the case of GPT-4 evaluations, the win rate of each method is determined by comparing them to the preferred completions present in the test set. Since no established SFT model exists for this scenario, our process begins with a pre-trained Pythia-2.8B. We use Preferred-FT to fine-tune a reference model on the preferred completions, ensuring that generated outputs remain close to the data distribution, and subsequently proceed with DPO training. For additional comparison, we include the Best of 128 Preferred-FT completions (our experiments indicate performance saturates beyond 128 completions for this benchmark; refer to Appendix Figure~\ref{fig:best-of-n}) as well as a 2-shot prompted variant of the original Pythia-2.8B model. We observe that DPO consistently matches or outperforms these alternatives at their optimal sampling temperatures. Furthermore, we evaluate an RLHF model based on PPO, trained on the Anthropic HH dataset \footnote{\url{https://huggingface.co/reciprocate/ppo_hh_pythia-6B}}, and made publicly available \footnote{\url{https://github.com/CarperAI/trlx/tree/main/examples/hh}}. However, no tested prompting or sampling configuration with PPO exceeds the baseline performance of Pythia-2.8B. Given both the results from TL;DR and the shared reward optimization objectives, we take Best of 128 as a practical indicator of PPO-like effectiveness. In summary, DPO stands out as the only method that is both computationally efficient and capable of surpassing the preferred completions benchmark in the Anthropic HH dataset, providing comparable or superior outcomes to the much costlier Best of 128 baseline. As illustrated in Figure~\ref{fig:dialogue-main}, DPO attains peak performance after relatively few iterations.

\subsection{Generalization to a new input distribution}

To assess the robustness of PPO and DPO under shifts in input distribution, we tested the policies learned in our Reddit TL;DR summarization task on a distinct dataset: the test split of CNN/DailyMail news articles \citep{nallapati-etal-2016-abstractive}. For this evaluation, we employed the same optimal sampling temperatures found for TL;DR (0 and 0.25). Table~\ref{tab:ood} reports these outcomes. We measured GPT-4's win rate by comparing its outputs to the ground-truth summaries using the identical GPT-4 (C) prompt as in the Reddit TL;DR setting, with ``forum post'' altered to ``news article''. On this new data distribution, DPO demonstrates a clear advantage over PPO, surpassing it by a wide margin. These findings provide early support that DPO-trained policies retain their generalization capacity, matching that of PPO, despite DPO not leveraging the supplementary unlabeled Reddit TL;DR prompts utilized in PPO.

\subsection{Validating GPT-4 judgments with human judgments}
\label{sec:human-judgments}
We carried out a human evaluation to cross-check the trustworthiness of GPT-4’s assessments by referencing results from the TL;DR summarization study and utilizing two distinct GPT-4 prompts. The \textbf{GPT-4 (S)} (simple) prompt asks which summary best covers the key information in the post, whereas the \textbf{GPT-4 (C)} (concise) prompt additionally requests a judgment of conciseness; this latter prompt is included because GPT-4, when following the \textbf{GPT-4 (S)} prompt, tends to favor summaries that are lengthier and more redundant than those preferred by human judges. Complete versions of the prompts appear in Appendix~\ref{app:prompts}. We conducted three comparison tests across methods to reflect a spectrum of sample quality: the highest-performing (DPO, temp. 0.25), lowest-performing (PPO, temp. 1.0), and a mid-range baseline (SFT, temp. 0.25); all three were benchmarked against PPO with greedy sampling (its optimal temperature). Results indicate that GPT-4’s preferences, for both prompts, align with human preferences to an extent comparable to inter-human agreement, implying GPT-4 serves as a robust stand-in for human assessment (due to a limited number of raters, only the DPO and PPO-1 comparisons received multiple independent human ratings). In summary, the \textbf{GPT-4 (C)} prompt tends to yield win rates that are more in line with human evaluation, which is why we adopt this prompt for key analyses in Section~\ref{sec:dpo-real-datasets}. For more information about the human study protocol, including screenshots of the rater interface and details on the human participant pool, please refer to Appendix~\ref{app:human-study}.

\section{Discussion}
Preference-based learning offers an effective and scalable approach for developing robust and aligned language models. In this work, we presented DPO, a straightforward method for leveraging human preferences in language model training that eliminates the need for reinforcement learning. Rather than reframing the task to fit into a typical RL context for compatibility with standard RL tools, DPO constructs a direct correspondence between language model policies and reward functions. This enables the direct optimization of language models in accordance with human feedback through a basic cross-entropy objective, foregoing both reinforcement learning and loss of expressive capability. DPO achieves comparable or superior outcomes to established RLHF approaches, including those utilizing PPO, all with minimal hyperparameter adjustment. This lowers the practical obstacles to preference-based language model training.

\textbf{Limitations \& Future Work.} Our findings prompt multiple directions for further investigation. One open question concerns the generalization abilities of DPO-trained policies outside the training distribution, relative to methods that use explicit reward modeling. Preliminary evidence indicates that DPO and PPO-based models show similar generalization, but systematic evaluation remains necessary. Another line of inquiry pertains to self-labeling: can DPO-trained policies facilitate effective learning from unlabeled prompts? Additionally, examining how reward over-optimization appears under direct preference optimization—and whether the modest drop in performance observed in Figure~\ref{fig:dialogue-main}-right reflects such effects—warrants further attention. While our experiments considered models of up to 6B parameters, future studies should assess the scalability of DPO for training next-generation models with much larger parameter counts. On the evaluation front, we observed that the win rates assigned by GPT-4 are sensitive to prompt formulation, suggesting that identifying optimal prompting strategies for automated evaluations is another area for future research. Beyond language modeling, DPO’s applicability to the training of generative models across diverse modalities presents exciting opportunities.